\documentclass[preprint,11pt,authoryear]{elsarticle}

\usepackage{amsmath,amssymb}
\usepackage{mathtools}
\usepackage[colorlinks=true,linkcolor=blue,citecolor=blue,urlcolor=blue]{hyperref}

\newtheorem{proposition}{Proposition}
\newtheorem{lemma}{Lemma}
\newtheorem{corollary}{Corollary}
\newdefinition{remark}{Remark}
\newproof{proof}{Proof}

\newcommand{\PnL}{\mathrm{PnL}}
\DeclarePairedDelimiter{\pos}{(}{)_{+}}

\setlength{\emergencystretch}{2em}

\journal{Economics Letters}

\begin{document}

\begin{frontmatter}

\title{Winning the large language capability battle\\ and losing the production economy}

\author{Peter Cotton}
\ead{peter.cotton@gsmc.ai}
\address{Global Strategic Minerals Corporation}

\begin{abstract}
The race to dominate machine intelligence is scored on capability benchmarks. This may be the wrong metric: in competitive equilibrium, work goes to whichever model creates the most value per token of compute, and we exhibit a closed form in which the most capable model's production compute share tends to zero. Task-level average PnL per token equals the shadow price of compute divided by the elasticity of task value with respect to effective compute, and ranks productivity exactly when that elasticity declines with crowding. No canonical task-invariant scalar quality exists: aggregate PnL-per-token rankings measure the economic niches a model occupies.
\end{abstract}

\begin{keyword}
division of labor \sep compute allocation \sep artificial intelligence \sep technological displacement \sep marginal product
\JEL D24 \sep D41 \sep O33
\end{keyword}

\end{frontmatter}

\section{Introduction}

The contest for leadership in large language models is treated as strategically decisive: governments restrict exports of the chips that train frontier models, laboratories commit capital at national-infrastructure scale, and progress is scored on capability benchmarks. Most public leaderboards rank capability, sometimes alongside price and latency; almost none ranks realized economic value net of compute expenditure. Beneath the race sits a premise: the most capable model captures the work, and with it the economic value of machine intelligence.

The premise fails in equilibrium. When compute is itinerant, free to move between prediction tasks and to run any available algorithm, and returns to compute on each task diminish with crowding, each subcontractible task goes to whichever model creates the most value per token of compute. Capability carries no allocation value beyond its contribution to that quotient: we exhibit a closed-form equilibrium in which a model that outscores every rival on every benchmark has production compute share tending to zero as the division of labor refines. The natural observable is profit-and-loss (PnL) per token in a prediction market run as a market scoring rule, which pays each entry its marginal contribution at the state it faced. Under sole supply, task-level average PnL per token equals the shadow price of compute divided by the elasticity of task value with respect to effective compute; at a common shadow price it ranks productivity when that elasticity declines with crowding. No comparable theorem covers an arbitrary lifetime quotient.

The mechanism is old. \citet{smith1776} tied the division of labour to the extent of the market; \citet{ricardo1817} showed the absolutely better party does not get all the work; \citet{beckermurphy1992} bound specialization by coordination costs; \citet{acemogluautor2011} model production as a task continuum with technologies assigned by comparative advantage; \citet{cotton2022} supplies the vision formalized here, of analytical tasks evolving into radically low-cost prediction supply chains. We import the skeleton, replace labor with compute, and add the prediction-specific feature: alpha decays with crowding.

\section{The economy}
\label{sec:economy}

Work is a unit mass $\Omega = [0,1]$ carrying Lebesgue measure $\mu$. At date $t$ it is partitioned into finitely many subcontractible tasks $\mathcal{P}_t = \{B_1, \ldots, B_{n_t}\}$, and partitions refine: every cell of $\mathcal{P}_t$ is a union of cells of $\mathcal{P}_{t+1}$. Finitely many models $m \in \mathcal{M}_t$ are available, with $\mathcal{M}_t \subseteq \mathcal{M}_{t+1}$. Each model $m$ has a feasibility set $A_m \subseteq \Omega$, a union of cells of $\mathcal{P}_t$ whenever $m \in \mathcal{M}_t$, and a productivity $\theta_m : A_m \to (0,\infty)$, constant on cells at date $t$; write $\theta_m(B)$ for its value on $B \subseteq A_m$ and require $q_{mB} = 0$ otherwise.

A \emph{token} is a standardized unit of compute cost, not a vendor text token, and PnL is expected trading gain under a fixed capital and risk protocol; without both normalizations PnL per token is not comparable across models.

Let $q_{mB} \ge 0$ be the density of tokens of model $m$ on task $B$. A token of model $m$ contributes $\theta_m(B)$ units of effective effort, so effort on $B$ is $z_B = \sum_m \theta_m(B)\, q_{mB}$, and the task pays $\mu(B)\,\phi(z_B)$, where $\phi$ is continuously differentiable on $[0,\infty)$ with $\phi(0) = 0$, $\phi' > 0$, $\phi'$ strictly decreasing, and $\phi'(z) \to 0$ as $z \to \infty$. The concavity says marginal alpha declines with crowding: prediction opportunities exhibit diminishing returns to compute. Assume $K > 0$, every cell has positive measure, and every cell is feasible for some model. Total value is
\begin{equation}
\label{eq:V}
V(q) = \sum_{B \in \mathcal{P}_t} \mu(B)\, \phi\!\Big( \sum_m \theta_m(B)\, q_{mB} \Big),
\qquad
\sum_B \mu(B) \sum_m q_{mB} = K.
\end{equation}

Compute is \emph{itinerant}: at zero switching cost a token may move to any task and run any available model. At the margin, a token of model $m$ on task $B$ earns $r_{mB} = \theta_m(B)\, \phi'(z_B)$. An equilibrium is an allocation and a return $\lambda$ with $r_{mB} = \lambda$ wherever $q_{mB} > 0$ and $r_{mB} \le \lambda$ elsewhere. In English: no token can raise its PnL by moving. Existence and efficiency are immediate: the routing game has concave potential $V$ over the compute simplex, and $\lambda$ is the shadow price of compute.

\begin{proposition}[Displacement]
\label{prop:displacement}
If $\theta_{m^*}(B) > \theta_m(B)$ for every $m \neq m^*$, then in any equilibrium supplying $B$, only $m^*$ runs on $B$.
\end{proposition}

\begin{proof}
If some $m \neq m^*$ had $q_{mB} > 0$ then
\begin{equation*}
\theta_m(B)\,\phi'(z_B) = \lambda
\qquad\text{and}\qquad
\theta_{m^*}(B)\,\phi'(z_B) > \lambda,
\end{equation*}
so a token improves by switching to $m^*$ on the same task. \qed
\end{proof}

Selection runs on the upper envelope $\theta^*(B) = \max_m \theta_m(B)$ over models feasible on $B$: an inferior technology receives zero execution share, not a small one. Displacement is this stark because value aggregates through the scalar $z_B$, making model outputs perfect substitutes after productivity scaling; models with complementary errors, whose ensemble improves on either alone, fall outside Proposition~\ref{prop:displacement}. Take $\phi(z) = \log(1+z)$. Active tasks satisfy $\theta^*(B) / (1 + \theta^*(B) q_B) = \lambda$, so
\begin{equation}
\label{eq:waterfill}
q_B = \pos*{\frac{1}{\lambda} - \frac{1}{\theta^*(B)}},
\qquad
K = \sum_B \mu(B) \pos*{\frac{1}{\lambda} - \frac{1}{\theta^*(B)}},
\end{equation}
the water-filling allocation; $\lambda$ and the aggregate intensities $q_B$ are unique, and the model-level split is unique when each active task's frontier model is unique.

\section{A more capable model exits}
\label{sec:exit}

Let $b_m$ be a model's success probability, $c_m$ its compute cost per unit of work, and $v$ any increasing conversion of accuracy into economic value, so that $\theta_m = v(b_m)/c_m$. The example needs only
\begin{equation}
b_G > b_S
\qquad\text{and, on the specialist's home tasks,}\qquad
\frac{v(b_G)}{c_G} < \frac{v(b_S)}{c_S}:
\end{equation}
more capable everywhere, less productive where the specialist is at home.

Take $v(b) = b$. A generalist $G$ succeeds with probability $0.99$ on every task at $c_G = 4$ tokens per unit of work, so $\theta_G = 0.2475$ everywhere. Specialist $S_k$ succeeds with probability $0.95$ on every task but is cheap only at home: with home tasks $I_k = (2^{-k}, 2^{-k+1}]$, its cost is $c_{S_k} = 1$ on $I_k$ and $10$ elsewhere, so $\theta_{S_k} = 0.95$ on $I_k$ and $0.095$ elsewhere. Every model is feasible everywhere, and the generalist is strictly more accurate than every specialist on every task. At date $t$ the partition distinguishes $I_1, \ldots, I_t$, leaving the residual $R_t = [0, 2^{-t}]$. By Proposition~\ref{prop:displacement} each $I_k$ goes entirely to $S_k$, the unique frontier there ($0.95 > 0.2475$); the generalist is the unique frontier only on the residual ($0.2475 > 0.095$), of measure $u_t = 2^{-t}$.

Compute share vanishes too, not just task measure. With $K = 4$, clearing \eqref{eq:waterfill} across the two active regions gives $\lambda_t = \left[4 + (1-u_t)/0.95 + u_t/0.2475\right]^{-1}$ and
\begin{equation}
Q_{G,t} = u_t \left( \frac{1}{\lambda_t} - \frac{1}{0.2475} \right)
= u_t \left( 1.0122 + 2.9878\, u_t \right) \longrightarrow 0.
\end{equation}
The more capable model exits production without ever losing a benchmark.

The example generalizes. Let $E_t = \{ x : \theta_G(x) \ge \max_{m \in \mathcal{M}_t,\, x \in A_m} \theta_m(x) \}$; the sets decrease, and by Proposition~\ref{prop:displacement} the generalist executes only on $E_t$. If specialists eventually dominate almost everywhere and equilibrium token intensities are bounded by $\bar q$, then $Q_{G,t} \le \bar q\, \mu(E_t) \to 0$. Conversely, at any fixed date, unique frontier status on a positive-measure, strictly active set guarantees positive compute share (a share bounded away from zero through time needs uniform measure and margin conditions); a null or inactive set supports none, and a tie leaves the share indeterminate. These are allocation results conditional on which technologies arrive, not predictions of arrival.

\section{What PnL per token measures}
\label{sec:measure}

Marginal PnL per token is not intrinsic to a model: it factorizes as $r_{mB} = \theta_m(B)\,\phi'(z_B)$, model times market state. Intrinsic, given the task, is $\theta_m(B)$; at matched task and congestion the market factor is common, so ratios of marginal PnL are ratios of $\theta$.

A market-scoring-rule mechanism, though not every market design, implements exactly this marginal-contribution accounting: each entry is paid its improvement over the standing forecast, and payments telescope \citep{hanson2007}. To identify the payments with the production technology, let $p(z)$ be the standing forecast after effective effort $z$ and \emph{define} $\phi(z) = \mathbb{E}[S(p(z), Y)] - \mathbb{E}[S(p(0), Y)]$, the score improvement effort has purchased, funded by a principal who procures the forecast (only these private returns are needed). An increment $dz = \theta\, dq$ earns $d\Pi = \theta\,\phi'(\theta q)\, dq$ for a sole supplier, the case Proposition~\ref{prop:displacement} delivers on strict frontiers (with interleaved suppliers, attribution is path dependent), and payments integrate to $\Pi(q) = \phi(\theta q)$. The form $\phi(z) = \log(1+z)$ is a parametrization, not a consequence of the scoring rule.

\begin{lemma}[PnL per token is an inverse elasticity]
\label{lem:ap}
Let $\varepsilon_\phi(z) = z\,\phi'(z)/\phi(z)$. On any active task with productivity $\theta$ at shadow price $\lambda$,
\begin{equation}
\frac{AP(\theta;\lambda)}{\lambda} = \frac{1}{\varepsilon_\phi(z)},
\qquad
\theta\,\phi'(z) = \lambda,
\end{equation}
and $z(\theta,\lambda)$ is increasing in $\theta$. Hence $AP$ ranks $\theta$ at a common $\lambda$ if and only if $\varepsilon_\phi$ decreases on the relevant range.
\end{lemma}

\begin{proof}
$AP = \phi(z)/q = \theta\,\phi(z)/z$ and $\lambda = \theta\,\phi'(z)$ give $AP/\lambda = \phi(z)/(z\,\phi'(z))$, and $z = (\phi')^{-1}(\lambda/\theta)$ increases in $\theta$. \qed
\end{proof}

\begin{corollary}[Log production ranks]
\label{cor:log}
For $\phi(z) = \log(1+z)$, $\varepsilon_\phi(z) = z/\big((1+z)\log(1+z)\big)$ is strictly decreasing, since $\log(1+z) < z$; hence $AP$ is strictly increasing in $\theta$, explicitly $AP = \lambda\, g(\theta/\lambda)$ with $g(r) = r\log r/(r-1)$ and $g'(r) \propto r - 1 - \log r > 0$.
\end{corollary}

Concavity alone does not deliver the ranking. At the Inada boundary, $\phi(z) = z^\alpha$ has constant elasticity, so $AP/\lambda = 1/\alpha$ for every $\theta$: no ranking information at all. Within the smooth class above, $\phi(z) = 0.1(1 - e^{-10z}) + (1 - e^{-0.1z})$ satisfies every stated condition yet its elasticity rises from about $0.33$ at $z = 0.4$ to $0.46$ at $z = 1$, reversing the ranking there. Whether the elasticity of alpha with respect to compute declines with crowding is the empirical condition on which the interpretation rests.

\begin{remark}[Elastic compute pins $\lambda = 1$]
\label{rem:elastic}
The $\lambda$-conditioning matters only when compute is capacity constrained. Elastic supply at price $c$ gives $\theta\,\phi'(\theta q) = c$, with a unique positive solution when $c < \theta\,\phi'(0)$ and $q = 0$ otherwise. A token being a standardized dollar of compute cost, normalize $c = 1$: then $\lambda = 1$ across unconstrained environments and $AP(\theta) = g(\theta)$ is directly comparable.
\end{remark}

In English: the margin is the same $\lambda$ for every survivor; the inframarginal rents are monotone in $\theta$. Mixing markets with different shadow prices can still invert the quotient, and the repair is to deflate, since $AP/\lambda$ depends on $\theta/\lambda$ alone: raw PnL per token is a nominal return, and the shadow value of compute is its deflator.

Administered deployment can invert the ranking within one environment, since it need not satisfy the common-$\lambda$ conditions: $\theta_A = 2$ forced to scale $q_A = 100$ averages $\log 201/100 \approx 0.053$ per token against $0.693$ for $\theta_B = 1$ at $q_B = 1$. A laboratory routing all traffic to its own generalist administers exactly this experiment.

Three statistics must be kept apart: the sole-supplier task-level average product, which obeys Lemma~\ref{lem:ap}; the incremental return $\Delta\PnL/\Delta T \approx \theta_m(B)\,\phi'(z_B)$ at matched task and congestion, the right statistic under randomized or interleaved deployment; and an arbitrary lifetime quotient, a deployment-weighted average of marginal returns to which no inverse-elasticity theorem applies.

One attempt at a realized-PnL ranking, \texttt{modelrank.net}, divides prediction-market PnL by vendor-metered tokens under randomized routing, with open positions marked to market. The quotient began as a heuristic for the survival-relevant statistic; the results above sharpen it. Its markets run on a logarithmic automated market maker,\footnote{Implementation details at \texttt{modelrank.net}.} the cost-function form of Hanson's rule, so each trade is paid its score increment and the telescoping identity applies literally; read through the model, the site is a randomized experiment estimating average incremental trading contribution over an assigned task distribution, randomization serving identification, not the itinerant equilibrium. The theory's statistic would denominate in dollar compute cost at matched task and state, deflated by the contemporaneous shadow price, under a common or suitably normalized task-response $\phi$. Nor is there a canonical task-invariant scalar underneath: a model serving many tasks reports a token-weighted portfolio of $\lambda\, g(\theta_m(B)/\lambda)$, so every aggregate ranking embeds a task distribution and deployment weighting, measuring the economic niches a model occupies. Capability benchmarks measure $b$ alone; equilibrium prices $b$ jointly with compute cost and task value through $\theta = v(b)/c$.

\section{Discussion}
\label{sec:discussion}

Capable large models need not disappear: their productivity advantage may move up the hierarchy, to decomposition, routing, verification, and the hard residual, whenever their value per token there sits on the frontier (in this model an application of Proposition~\ref{prop:displacement}, since decomposition is free and exogenous). For strategy, the question is not who owns the most capable model but who owns the productivity frontier on an economically material part of the limiting division of labor. PnL per token is the intensive margin; production share also needs the extensive margin, the economic mass of the niches occupied, and $\theta = 100$ on a niche of measure $10^{-9}$ is spectacular PnL per token with negligible share. Commercial survival needs more still: an active task, $\theta > \lambda$, yields rent $R(\theta,\lambda) = \log(\theta/\lambda) - 1 + \lambda/\theta$ per unit of task mass above the shadow opportunity cost of its compute, zero if inactive ($\lambda = 1$ and $R(\theta) = \log\theta - 1 + 1/\theta$ in cash under Remark~\ref{rem:elastic}), and it is total rent across owned tasks that must cover fixed cost. The chain runs: task productivity, PnL per token, allocation, total rents, commercial survival. This note proves the middle of the chain; the entry model closing it is left to future work.

\section*{Declaration of generative AI and AI-assisted technologies in the manuscript preparation process}

During the preparation of this work, the author used Claude (Anthropic) for drafting, typesetting, and the checking of derivations. The author reviewed and edited the output as needed and takes full responsibility for the content of the published article.

\bibliographystyle{elsarticle-harv}

\small
\begin{thebibliography}{4}
\providecommand{\natexlab}[1]{#1}

\bibitem[Acemoglu and Autor(2011)]{acemogluautor2011}
Acemoglu, D., Autor, D., 2011. Skills, tasks and technologies: Implications for
employment and earnings, in: Ashenfelter, O., Card, D. (Eds.), Handbook of
Labor Economics, Vol. 4B. Elsevier, pp. 1043--1171.

\bibitem[Becker and Murphy(1992)]{beckermurphy1992}
Becker, G.S., Murphy, K.M., 1992. The division of labor, coordination costs,
and knowledge. Quarterly Journal of Economics 107, 1137--1160.

\bibitem[Cotton(2022)]{cotton2022}
Cotton, P., 2022. Microprediction: Building an Open AI Network. MIT Press,
Cambridge, MA.

\bibitem[Hanson(2007)]{hanson2007}
Hanson, R., 2007. Logarithmic market scoring rules for modular combinatorial
information aggregation. Journal of Prediction Markets 1, 3--15.

\bibitem[Ricardo(1817)]{ricardo1817}
Ricardo, D., 1817. On the Principles of Political Economy and Taxation.
John Murray, London.

\bibitem[Smith(1776)]{smith1776}
Smith, A., 1776. An Inquiry into the Nature and Causes of the Wealth of
Nations. W. Strahan and T. Cadell, London.

\end{thebibliography}

\end{document}
